% appendices/appendix_commands.tex
\chapter*{附录B：dmctl 命令参考}
\addcontentsline{toc}{chapter}{附录B：dmctl 命令参考}

\section*{初始化与原始 I/O 命令}

\begin{longtable}{|p{0.15\textwidth}|p{0.35\textwidth}|p{0.40\textwidth}|}
\hline
\textbf{命令} & \textbf{签名与示例} & \textbf{实现位置} \\
\hline
\endhead
init-image & \textbf{签名}：init-image <path> <blocks> <block\_size> & main.rs:32 \\
 & \textbf{示例}：dmctl init-image disk.img 1024 4096 & \\
 & \textbf{说明}：创建零填充文件用作虚拟块设备 & \\
\hline
raw-read-hex & \textbf{签名}：raw-read-hex <path> <block\_size> <block\_id> & main.rs:65 \\
 & \textbf{示例}：dmctl raw-read-hex disk.img 4096 0 & \\
 & \textbf{说明}：读取块以十六进制打印（绕过映射器） & \\
\hline
tamper & \textbf{签名}：tamper <path> <block\_size> <block\_id> <offset> <hex\_byte> & main.rs:133-137 \\
 & \textbf{示例}：dmctl tamper disk.img 4096 0 10 ff & \\
 & \textbf{说明}：修改块内字节（测试 verity） & \\
\hline
\end{longtable}

\section*{加密命令（dm-crypt）}

\begin{longtable}{|p{0.15\textwidth}|p{0.35\textwidth}|p{0.40\textwidth}|}
\hline
\textbf{命令} & \textbf{签名与示例} & \textbf{实现位置} \\
\hline
\endhead
crypt-write & \textbf{签名}：crypt-write <path> <block\_size> <key> <block\_id> <data\_string> & main.rs:41 \\
 & \textbf{示例}：dmctl crypt-write disk.img 4096 mysecretkey 0 "Hello" & \\
\hline
crypt-read & \textbf{签名}：crypt-read <path> <block\_size> <key> <block\_id> & main.rs \\
 & \textbf{示例}：dmctl crypt-read disk.img 4096 mysecretkey 0 & \\
\hline
crypt-table-fill & \textbf{签名}：crypt-table-fill <path> <block\_size> <table\_str> <start> <byte> <count> & main.rs:72-77 \\
 & \textbf{示例}：dmctl crypt-table-fill disk.img 4096 "aes-xts-plain64 <key> 0 /dev/loop0 0" 0 aa 10 & \\
\hline
\end{longtable}

\section*{完整性命令（dm-verity）}

\begin{longtable}{|p{0.15\textwidth}|p{0.35\textwidth}|p{0.40\textwidth}|}
\hline
\textbf{命令} & \textbf{签名与示例} & \textbf{实现位置} \\
\hline
\endhead
verity-root & \textbf{签名}：verity-root <path> <block\_size> & main.rs:91 \\
 & \textbf{示例}：dmctl verity-root disk.img 4096 & \\
\hline
verity-verify & \textbf{签名}：verity-verify <path> <block\_size> <root\_hash\_hex> & main.rs:104-106 \\
 & \textbf{示例}：dmctl verity-verify disk.img 4096 a1b2... & \\
\hline
verity-format-dm & \textbf{签名}：verity-format-dm <path> <data\_blocks> & main.rs:113 \\
 & \textbf{示例}：dmctl verity-format-dm disk.img 1024 & \\
 & \textbf{说明}：为星绽内核格式化镜像 & \\
\hline
\end{longtable}

\section*{演示与基准测试命令}

\begin{longtable}{|p{0.15\textwidth}|p{0.35\textwidth}|p{0.40\textwidth}|}
\hline
\textbf{命令} & \textbf{签名与示例} & \textbf{实现位置} \\
\hline
\endhead
demo-crypt & \textbf{签名}：demo-crypt & main.rs:140-141 \\
 & \textbf{示例}：dmctl demo-crypt & \\
\hline
demo-verity & \textbf{签名}：demo-verity & main.rs:140-141 \\
 & \textbf{示例}：dmctl demo-verity & \\
\hline
benchmark & \textbf{签名}：benchmark [blocks] [iterations] & main.rs:142-146 \\
 & \textbf{示例}：dmctl benchmark 4096 5 & \\
\hline
\end{longtable}